% --- Archon physics formalization source begin ---
% archon:physics
% archon:covers PhyXMiniProblems/problem_phyx_mini_0601.lean
% archon:source-report reports/phyx_mini/problem_phyx_mini_0601.source.json

\chapter{Physics problem phyx\_mini\_0601}
\label{ch:PhyXMiniProblems_problem_phyx_mini_0601}

\paragraph{Problem source.}
\#\# Physical scenario

The Scattering Matrix. The theory of scattering generalizes in a pretty obvious way to arbitrary localized potentials. To the left (Region I), \$V(x) = 0\$, so \$\$\textbackslash{}psi(x) = Ae\textasciicircum{}\{ikx\} + Be\textasciicircum{}\{-ikx\}, \textbackslash{}quad \textbackslash{}text\{where \} k = \textbackslash{}sqrt\{\textbackslash{}frac\{2mE\}\{\textbackslash{}hbar\}\}.\$\$ To the right (Region III), \$V(x)\$ is again zero, so \$\$\textbackslash{}psi(x) = Fe\textasciicircum{}\{ikx\} + Ge\textasciicircum{}\{-ikx\}.\$\$ In between (Region II), of course, I can't tell you what \$\textbackslash{}psi\$ is until you specify the potential, but because the Schrödinger equation is a linear, second-order differential equation, the general solution has got to be of the form \$\$\textbackslash{}psi(x) = Cf(x) + Dg(x),\$\$ where \$f(x)\$ and \$g(x)\$ are two linearly independent particular solutions. There will be four boundary conditions (two joining Regions I and II, and two joining Regions II and III). Two of these can be used to eliminate \$C\$ and \$D\$, and the other two can be "solved" for \$B\$ and \$F\$ in terms of \$A\$ and \$G\$: \$\$B = S\_\{11\}A + S\_\{12\}G, \textbackslash{}quad F = S\_\{21\}A + S\_\{22\}G.\$\$ The four coefficients \$S\_\{ij\}\$, which depend on \$k\$ (and hence on \$E\$), constitute a \$2 \textbackslash{}times 2\$ matrix \$\textbackslash{}mathbf\{S\}\$, called the scattering matrix (or \$S\$-matrix, for short). The \$S\$-matrix tells you the outgoing amplitudes (\$B\$ and \$F\$) in terms of the incoming amplitudes (\$A\$ and \$G\$): \$\$egin\{pmatrix\} B \textbackslash{} F \textbackslash{}end\{pmatrix\} = egin\{pmatrix\} S\_\{11\} \& S\_\{12\} \textbackslash{} S\_\{21\} \& S\_\{22\} \textbackslash{}end\{pmatrix\} egin\{pmatrix\} A \textbackslash{} G \textbackslash{}end\{pmatrix\}.\$\$ In the typical case of scattering from the left, \$G = 0\$, so the reflection and transmission coefficients are \$\$R\_L = \textbackslash{}left| \textbackslash{}frac\{B\}\{A\} 
ight|\_\{G=0\}\textasciicircum{}2 = |S\_\{11\}|\textasciicircum{}2, \textbackslash{}quad T\_L = \textbackslash{}left| \textbackslash{}frac\{F\}\{A\} 
ight|\_\{G=0\}\textasciicircum{}2 = |S\_\{21\}|\textasciicircum{}2.\$\$ For scattering from the right, \$A = 0\$, and \$\$R\_R = \textbackslash{}left| \textbackslash{}frac\{F\}\{G\} 
ight|\_\{A=0\}\textasciicircum{}2 = |S\_\{22\}|\textasciicircum{}2, \textbackslash{}quad T\_R = \textbackslash{}left| \textbackslash{}frac\{B\}\{G\} 
ight|\_\{A=0\}\textasciicircum{}2 = |S\_\{12\}|\textasciicircum{}2.\$\$

\#\# Question

Construct the \$S\$-matrix for scattering from a delta-function well

\#\# Answer choices

- A: A: \textbackslash{}left| \textbackslash{}frac\{1\}\{S\_\{22\}\} 
ight|\textasciicircum{}2

- B: B: \textbackslash{}left| \textbackslash{}frac\{M\_\{12\}\}\{M\_\{22\}\} 
ight|\textasciicircum{}2

- C: C: \textbackslash{}frac\{1\}\{|M\_\{22\}|\textasciicircum{}2\}

- D: D: |\textbackslash{}det(M)|\textasciicircum{}2

\#\# Figure caption (auxiliary; use the image as primary evidence)

The image is a diagram illustrating a potential barrier in quantum mechanics. It is divided into three regions: Region I, Region II, and Region III. Here’s a breakdown of the content:

1. **Axes**:
   - The horizontal axis is labeled \textbackslash{}( x \textbackslash{}), presumably representing position.
   - The vertical axis represents potential energy \textbackslash{}( V(x) \textbackslash{}).

2. **Potential Curve**:
   - There is a curve showing variations in potential energy, \textbackslash{}( V(x) \textbackslash{}), within Region II. 
   - The curve starts at a baseline, rises to a peak, dips down below the baseline, and then returns to near the baseline as it moves from left to right.

3. **Regions**:
   - **Region I** is on the far left, where the potential is constant.
   - **Region II** is in the center, where the potential varies.
   - **Region III** is on the far right, similar to Region I with constant potential.

4. **Wave Functions**:
   - In Region I, the wave functions are depicted as \textbackslash{}( Ae\textasciicircum{}\{ikx\} \textbackslash{}) (rightward arrow) and \textbackslash{}( Be\textasciicircum{}\{-ikx\} \textbackslash{}) (leftward arrow), indicating incident and reflected waves.
   - In Region III, the wave functions are \textbackslash{}( Fe\textasciicircum{}\{ikx\} \textbackslash{}) (rightward arrow) and \textbackslash{}( Ge\textasciicircum{}\{-ikx\} \textbackslash{}) (leftward arrow), indicating transmitted and reflected waves.

5. **Arrows**:
   - Arrows above or below the wave functions in Regions I and III indicate the direction of wave propagation (rightward for \textbackslash{}( Ae\textasciicircum{}\{ikx\} \textbackslash{}) and \textbackslash{}( Fe\textasciicircum{}\{ikx\} \textbackslash{}), leftward for \textbackslash{}( Be\textasciicircum{}\{-ikx\} \textbackslash{}) and \textbackslash{}( Ge\textasciicircum{}\{-ikx\} \textbackslash{})).

Overall, the diagram visually represents the interaction of wave functions with a potential barrier, relevant to quantum tunneling or scattering.

\#\# Dataset metadata

Quantum Phenomena; Physical Model Grounding Reasoning, Numerical Reasoning

\paragraph{Recorded answer/context.}
C: C: \textbackslash{}frac\{1\}\{|M\_\{22\}|\textasciicircum{}2\}

\paragraph{Figure/image path.}
/root/proposal\_for\_physic/science-mango/phyx\_mini\_run/phyx\_data/test\_image/601.png

\paragraph{Formalization target.}
create a compiling Lean file with sorry bodies at `PhyXMiniProblems/problem\_phyx\_mini\_0601.lean`. The Lean declarations must preserve the physical quantities, dimensions or dimensional roles, figure labels, governing-law hypotheses, and final relation expressed by this problem.
Use Mathlib/Physlib names found through LeanExplore where available. If a physics API is missing, introduce faithful local abstractions rather than scalar placeholder aliases.

\begin{theorem}[Physics formalization target]
\label{thm:physics:phyx_mini_0601:target}
\lean{PhyXMiniProblems.ProblemPhyXMini0601.problem_phyx_mini_0601}
\uses{def:physics:phyx-mini-0601:phyxminiproblems-problemphyxmini0601-deltawellscatteringsetup, def:physics:phyx-mini-0601:phyxminiproblems-problemphyxmini0601-matchesprimaryscatteringfigure, def:physics:phyx-mini-0601:phyxminiproblems-problemphyxmini0601-matchesattractivedeltawellscenario, def:physics:phyx-mini-0601:phyxminiproblems-problemphyxmini0601-hasphysicaldeltawellparameters, def:physics:phyx-mini-0601:phyxminiproblems-problemphyxmini0601-satisfiesfreeparticledispersionlaw, def:physics:phyx-mini-0601:phyxminiproblems-problemphyxmini0601-satisfiesasymptoticwaveforms, def:physics:phyx-mini-0601:phyxminiproblems-problemphyxmini0601-satisfiesregioniisolutiondescription, def:physics:phyx-mini-0601:phyxminiproblems-problemphyxmini0601-deltawellcoupling, def:physics:phyx-mini-0601:phyxminiproblems-problemphyxmini0601-satisfiesdeltawellboundarylaws, def:physics:phyx-mini-0601:phyxminiproblems-problemphyxmini0601-satisfiesscatteringmatrixrole, def:physics:phyx-mini-0601:phyxminiproblems-problemphyxmini0601-satisfiestransfermatrixlaws, def:physics:phyx-mini-0601:phyxminiproblems-problemphyxmini0601-isleftincidenttrial, def:physics:phyx-mini-0601:phyxminiproblems-problemphyxmini0601-satisfiesscatteringobservables, def:physics:phyx-mini-0601:phyxminiproblems-problemphyxmini0601-matchesanswerchoice, def:physics:phyx-mini-0601:phyxminiproblems-problemphyxmini0601-recordeddatasetanswer}
The assigned autoformalize agent should translate this physics problem into Lean declarations in the covered file, with theorem and lemma proof bodies written as `by sorry`.
\end{theorem}
\begin{proof}
This is an autoformalization task, not a proof task. Produce statements that can later be proved by the physics prover without weakening the physical model.
\end{proof}

% --- Archon declaration topology begin ---
\section{Declaration topology}
\begin{definition}[Mass Quantity]
  \label{def:physics:phyx-mini-0601:phyxminiproblems-problemphyxmini0601-massquantity}
  \lean{PhyXMiniProblems.ProblemPhyXMini0601.MassQuantity}
  A unit-independent particle mass, of physical dimension M
\end{definition}
\begin{proof}
  This is the declared model definition of Mass Quantity.
\end{proof}
\begin{definition}[Wave Number Quantity]
  \label{def:physics:phyx-mini-0601:phyxminiproblems-problemphyxmini0601-wavenumberquantity}
  \lean{PhyXMiniProblems.ProblemPhyXMini0601.WaveNumberQuantity}
  A unit-independent one-dimensional wave number, of dimension L⁻¹
\end{definition}
\begin{proof}
  This is the declared model definition of Wave Number Quantity.
\end{proof}
\begin{definition}[Delta Well Strength Quantity]
  \label{def:physics:phyx-mini-0601:phyxminiproblems-problemphyxmini0601-deltawellstrengthquantity}
  \lean{PhyXMiniProblems.ProblemPhyXMini0601.DeltaWellStrengthQuantity}
  The positive strength α in the attractive potential V(x) = -α δ(x-x₀). It has dimensions energy times length, M L³ T⁻²
\end{definition}
\begin{proof}
  This is the declared model definition of Delta Well Strength Quantity.
\end{proof}
\begin{definition}[mass In Kilograms]
  \label{def:physics:phyx-mini-0601:phyxminiproblems-problemphyxmini0601-massinkilograms}
  \lean{PhyXMiniProblems.ProblemPhyXMini0601.massInKilograms}
  \uses{def:physics:phyx-mini-0601:phyxminiproblems-problemphyxmini0601-massquantity}
  Coherent-SI mass readout, in kilograms
\end{definition}
\begin{proof}
  This is the declared model definition of mass In Kilograms.
\end{proof}
\begin{definition}[energy In Joules]
  \label{def:physics:phyx-mini-0601:phyxminiproblems-problemphyxmini0601-energyinjoules}
  \lean{PhyXMiniProblems.ProblemPhyXMini0601.energyInJoules}
  Coherent-SI energy readout, in joules
\end{definition}
\begin{proof}
  This is the declared model definition of energy In Joules.
\end{proof}
\begin{definition}[wave Number In Inverse Meters]
  \label{def:physics:phyx-mini-0601:phyxminiproblems-problemphyxmini0601-wavenumberininversemeters}
  \lean{PhyXMiniProblems.ProblemPhyXMini0601.waveNumberInInverseMeters}
  \uses{def:physics:phyx-mini-0601:phyxminiproblems-problemphyxmini0601-wavenumberquantity}
  Coherent-SI wave-number readout, in inverse metres
\end{definition}
\begin{proof}
  This is the declared model definition of wave Number In Inverse Meters.
\end{proof}
\begin{definition}[delta Well Strength In Joule Meters]
  \label{def:physics:phyx-mini-0601:phyxminiproblems-problemphyxmini0601-deltawellstrengthinjoulemeters}
  \lean{PhyXMiniProblems.ProblemPhyXMini0601.deltaWellStrengthInJouleMeters}
  \uses{def:physics:phyx-mini-0601:phyxminiproblems-problemphyxmini0601-deltawellstrengthquantity}
  Coherent-SI delta-well strength readout, in joule metres
\end{definition}
\begin{proof}
  This is the declared model definition of delta Well Strength In Joule Meters.
\end{proof}
\begin{definition}[reduced Planck Constant In Joule Seconds]
  \label{def:physics:phyx-mini-0601:phyxminiproblems-problemphyxmini0601-reducedplanckconstantinjouleseconds}
  \lean{PhyXMiniProblems.ProblemPhyXMini0601.reducedPlanckConstantInJouleSeconds}
  Physlib's reduced Planck constant, read in joule seconds
\end{definition}
\begin{proof}
  This is the declared model definition of reduced Planck Constant In Joule Seconds.
\end{proof}
\begin{definition}[Scattering Region]
  \label{def:physics:phyx-mini-0601:phyxminiproblems-problemphyxmini0601-scatteringregion}
  \lean{PhyXMiniProblems.ProblemPhyXMini0601.ScatteringRegion}
  The three spatial regions printed under the potential diagram
\end{definition}
\begin{proof}
  This is the declared model definition of Scattering Region.
\end{proof}
\begin{definition}[Propagation Direction]
  \label{def:physics:phyx-mini-0601:phyxminiproblems-problemphyxmini0601-propagationdirection}
  \lean{PhyXMiniProblems.ProblemPhyXMini0601.PropagationDirection}
  Horizontal direction of propagation of a plane-wave component
\end{definition}
\begin{proof}
  This is the declared model definition of Propagation Direction.
\end{proof}
\begin{definition}[Asymptotic Wave Component]
  \label{def:physics:phyx-mini-0601:phyxminiproblems-problemphyxmini0601-asymptoticwavecomponent}
  \lean{PhyXMiniProblems.ProblemPhyXMini0601.AsymptoticWaveComponent}
  The four amplitude labels attached to arrows in the primary image
\end{definition}
\begin{proof}
  This is the declared model definition of Asymptotic Wave Component.
\end{proof}
\begin{definition}[Figure Axis]
  \label{def:physics:phyx-mini-0601:phyxminiproblems-problemphyxmini0601-figureaxis}
  \lean{PhyXMiniProblems.ProblemPhyXMini0601.FigureAxis}
  The two axes appearing in the potential diagram
\end{definition}
\begin{proof}
  This is the declared model definition of Figure Axis.
\end{proof}
\begin{definition}[Axis Symbol]
  \label{def:physics:phyx-mini-0601:phyxminiproblems-problemphyxmini0601-axissymbol}
  \lean{PhyXMiniProblems.ProblemPhyXMini0601.AxisSymbol}
  Literal axis symbols printed in the image
\end{definition}
\begin{proof}
  This is the declared model definition of Axis Symbol.
\end{proof}
\begin{definition}[Scattering Figure]
  \label{def:physics:phyx-mini-0601:phyxminiproblems-problemphyxmini0601-scatteringfigure}
  \lean{PhyXMiniProblems.ProblemPhyXMini0601.ScatteringFigure}
  \uses{def:physics:phyx-mini-0601:phyxminiproblems-problemphyxmini0601-scatteringregion, def:physics:phyx-mini-0601:phyxminiproblems-problemphyxmini0601-propagationdirection, def:physics:phyx-mini-0601:phyxminiproblems-problemphyxmini0601-asymptoticwavecomponent, def:physics:phyx-mini-0601:phyxminiproblems-problemphyxmini0601-figureaxis, def:physics:phyx-mini-0601:phyxminiproblems-problemphyxmini0601-axissymbol}
  Raster-visible qualitative data from image 601. The image is a generic localized-potential diagram rather than a literal drawing of a Dirac delta; the delta-well specialization is therefore kept in a separate scenario law
\end{definition}
\begin{proof}
  This is the declared model definition of Scattering Figure.
\end{proof}
\begin{definition}[Delta Well Scattering Setup]
  \label{def:physics:phyx-mini-0601:phyxminiproblems-problemphyxmini0601-deltawellscatteringsetup}
  \lean{PhyXMiniProblems.ProblemPhyXMini0601.DeltaWellScatteringSetup}
  \uses{def:physics:phyx-mini-0601:phyxminiproblems-problemphyxmini0601-massquantity, def:physics:phyx-mini-0601:phyxminiproblems-problemphyxmini0601-wavenumberquantity, def:physics:phyx-mini-0601:phyxminiproblems-problemphyxmini0601-deltawellstrengthquantity, def:physics:phyx-mini-0601:phyxminiproblems-problemphyxmini0601-scatteringfigure}
  All quantities required by the prose, the delta-well specialization, and the four answer choices. A, G are incoming amplitudes; B, F are outgoing. The distribution field is an SI-coordinate representation of the potential
\end{definition}
\begin{proof}
  This is the declared model definition of Delta Well Scattering Setup.
\end{proof}
\begin{definition}[region IPlane Wave]
  \label{def:physics:phyx-mini-0601:phyxminiproblems-problemphyxmini0601-regioniplanewave}
  \lean{PhyXMiniProblems.ProblemPhyXMini0601.regionIPlaneWave}
  \uses{def:physics:phyx-mini-0601:phyxminiproblems-problemphyxmini0601-wavenumberininversemeters, def:physics:phyx-mini-0601:phyxminiproblems-problemphyxmini0601-deltawellscatteringsetup}
  The Region-I plane-wave superposition Aeⁱᵏˣ + Be⁻ⁱᵏˣ
\end{definition}
\begin{proof}
  This is the declared model definition of region IPlane Wave.
\end{proof}
\begin{definition}[region IIIPlane Wave]
  \label{def:physics:phyx-mini-0601:phyxminiproblems-problemphyxmini0601-regioniiiplanewave}
  \lean{PhyXMiniProblems.ProblemPhyXMini0601.regionIIIPlaneWave}
  \uses{def:physics:phyx-mini-0601:phyxminiproblems-problemphyxmini0601-wavenumberininversemeters, def:physics:phyx-mini-0601:phyxminiproblems-problemphyxmini0601-deltawellscatteringsetup}
  The Region-III plane-wave superposition Feⁱᵏˣ + Ge⁻ⁱᵏˣ
\end{definition}
\begin{proof}
  This is the declared model definition of region IIIPlane Wave.
\end{proof}
\begin{definition}[Are Linearly Independent Complex Functions]
  \label{def:physics:phyx-mini-0601:phyxminiproblems-problemphyxmini0601-arelinearlyindependentcomplexfunctions}
  \lean{PhyXMiniProblems.ProblemPhyXMini0601.AreLinearlyIndependentComplexFunctions}
  Elementary linear independence of the two particular complex-valued solutions used in Region II
\end{definition}
\begin{proof}
  This is the declared model definition of Are Linearly Independent Complex Functions.
\end{proof}
\begin{definition}[Matches Primary Scattering Figure]
  \label{def:physics:phyx-mini-0601:phyxminiproblems-problemphyxmini0601-matchesprimaryscatteringfigure}
  \lean{PhyXMiniProblems.ProblemPhyXMini0601.MatchesPrimaryScatteringFigure}
  \uses{def:physics:phyx-mini-0601:phyxminiproblems-problemphyxmini0601-deltawellscatteringsetup}
  Exact qualitative labels and arrow directions visible in image 601
\end{definition}
\begin{proof}
  This is the declared model definition of Matches Primary Scattering Figure.
\end{proof}
\begin{definition}[Matches Attractive Delta Well Scenario]
  \label{def:physics:phyx-mini-0601:phyxminiproblems-problemphyxmini0601-matchesattractivedeltawellscenario}
  \lean{PhyXMiniProblems.ProblemPhyXMini0601.MatchesAttractiveDeltaWellScenario}
  \uses{def:physics:phyx-mini-0601:phyxminiproblems-problemphyxmini0601-deltawellstrengthinjoulemeters, def:physics:phyx-mini-0601:phyxminiproblems-problemphyxmini0601-deltawellscatteringsetup}
  The question's attractive delta-function well, positioned at the origin. A Dirac distribution is used instead of replacing the potential with a scalar
\end{definition}
\begin{proof}
  This is the declared model definition of Matches Attractive Delta Well Scenario.
\end{proof}
\begin{definition}[Has Physical Delta Well Parameters]
  \label{def:physics:phyx-mini-0601:phyxminiproblems-problemphyxmini0601-hasphysicaldeltawellparameters}
  \lean{PhyXMiniProblems.ProblemPhyXMini0601.HasPhysicalDeltaWellParameters}
  \uses{def:physics:phyx-mini-0601:phyxminiproblems-problemphyxmini0601-massinkilograms, def:physics:phyx-mini-0601:phyxminiproblems-problemphyxmini0601-energyinjoules, def:physics:phyx-mini-0601:phyxminiproblems-problemphyxmini0601-wavenumberininversemeters, def:physics:phyx-mini-0601:phyxminiproblems-problemphyxmini0601-deltawellstrengthinjoulemeters, def:physics:phyx-mini-0601:phyxminiproblems-problemphyxmini0601-deltawellscatteringsetup}
  Positive mass, energy, wave number, and nonzero attractive-well strength
\end{definition}
\begin{proof}
  This is the declared model definition of Has Physical Delta Well Parameters.
\end{proof}
\begin{definition}[source Printed Dispersion Expression SI]
  \label{def:physics:phyx-mini-0601:phyxminiproblems-problemphyxmini0601-sourceprinteddispersionexpressionsi}
  \lean{PhyXMiniProblems.ProblemPhyXMini0601.sourcePrintedDispersionExpressionSI}
  \uses{def:physics:phyx-mini-0601:phyxminiproblems-problemphyxmini0601-massinkilograms, def:physics:phyx-mini-0601:phyxminiproblems-problemphyxmini0601-energyinjoules, def:physics:phyx-mini-0601:phyxminiproblems-problemphyxmini0601-reducedplanckconstantinjouleseconds, def:physics:phyx-mini-0601:phyxminiproblems-problemphyxmini0601-deltawellscatteringsetup}
  The scalar expression printed on the right-hand side of the source's formula k = sqrt (2 m E / ℏ). It is retained as source metadata, not asserted as a governing law, because it is dimensionally inconsistent
\end{definition}
\begin{proof}
  This is the declared model definition of source Printed Dispersion Expression SI.
\end{proof}
\begin{definition}[Satisfies Free Particle Dispersion Law]
  \label{def:physics:phyx-mini-0601:phyxminiproblems-problemphyxmini0601-satisfiesfreeparticledispersionlaw}
  \lean{PhyXMiniProblems.ProblemPhyXMini0601.SatisfiesFreeParticleDispersionLaw}
  \uses{def:physics:phyx-mini-0601:phyxminiproblems-problemphyxmini0601-massinkilograms, def:physics:phyx-mini-0601:phyxminiproblems-problemphyxmini0601-energyinjoules, def:physics:phyx-mini-0601:phyxminiproblems-problemphyxmini0601-wavenumberininversemeters, def:physics:phyx-mini-0601:phyxminiproblems-problemphyxmini0601-reducedplanckconstantinjouleseconds, def:physics:phyx-mini-0601:phyxminiproblems-problemphyxmini0601-deltawellscatteringsetup}
  The physically standard free-particle dispersion relation k² = 2 m E / ℏ², stated through coherent-SI readouts. Unlike the printed source expression above, both sides have dimensions of inverse metres squared
\end{definition}
\begin{proof}
  This is the declared model definition of Satisfies Free Particle Dispersion Law.
\end{proof}
\begin{definition}[Satisfies Asymptotic Wave Forms]
  \label{def:physics:phyx-mini-0601:phyxminiproblems-problemphyxmini0601-satisfiesasymptoticwaveforms}
  \lean{PhyXMiniProblems.ProblemPhyXMini0601.SatisfiesAsymptoticWaveForms}
  \uses{def:physics:phyx-mini-0601:phyxminiproblems-problemphyxmini0601-deltawellscatteringsetup, def:physics:phyx-mini-0601:phyxminiproblems-problemphyxmini0601-regioniplanewave, def:physics:phyx-mini-0601:phyxminiproblems-problemphyxmini0601-regioniiiplanewave}
  The asymptotically free Region-I and Region-III wave forms in the prose
\end{definition}
\begin{proof}
  This is the declared model definition of Satisfies Asymptotic Wave Forms.
\end{proof}
\begin{definition}[Satisfies Region IISolution Description]
  \label{def:physics:phyx-mini-0601:phyxminiproblems-problemphyxmini0601-satisfiesregioniisolutiondescription}
  \lean{PhyXMiniProblems.ProblemPhyXMini0601.SatisfiesRegionIISolutionDescription}
  \uses{def:physics:phyx-mini-0601:phyxminiproblems-problemphyxmini0601-deltawellscatteringsetup, def:physics:phyx-mini-0601:phyxminiproblems-problemphyxmini0601-arelinearlyindependentcomplexfunctions}
  The source's Region-II statement ψ = C f + D g, with f and g linearly independent particular solutions of the linear second-order equation
\end{definition}
\begin{proof}
  This is the declared model definition of Satisfies Region IISolution Description.
\end{proof}
\begin{definition}[delta Well Coupling]
  \label{def:physics:phyx-mini-0601:phyxminiproblems-problemphyxmini0601-deltawellcoupling}
  \lean{PhyXMiniProblems.ProblemPhyXMini0601.deltaWellCoupling}
  \uses{def:physics:phyx-mini-0601:phyxminiproblems-problemphyxmini0601-massinkilograms, def:physics:phyx-mini-0601:phyxminiproblems-problemphyxmini0601-wavenumberininversemeters, def:physics:phyx-mini-0601:phyxminiproblems-problemphyxmini0601-deltawellstrengthinjoulemeters, def:physics:phyx-mini-0601:phyxminiproblems-problemphyxmini0601-reducedplanckconstantinjouleseconds, def:physics:phyx-mini-0601:phyxminiproblems-problemphyxmini0601-deltawellscatteringsetup}
  Dimensionless coupling η = m α / (ℏ² k) for an attractive delta well
\end{definition}
\begin{proof}
  This is the declared model definition of delta Well Coupling.
\end{proof}
\begin{definition}[Satisfies Delta Well Boundary Laws]
  \label{def:physics:phyx-mini-0601:phyxminiproblems-problemphyxmini0601-satisfiesdeltawellboundarylaws}
  \lean{PhyXMiniProblems.ProblemPhyXMini0601.SatisfiesDeltaWellBoundaryLaws}
  \uses{def:physics:phyx-mini-0601:phyxminiproblems-problemphyxmini0601-massinkilograms, def:physics:phyx-mini-0601:phyxminiproblems-problemphyxmini0601-wavenumberininversemeters, def:physics:phyx-mini-0601:phyxminiproblems-problemphyxmini0601-deltawellstrengthinjoulemeters, def:physics:phyx-mini-0601:phyxminiproblems-problemphyxmini0601-reducedplanckconstantinjouleseconds, def:physics:phyx-mini-0601:phyxminiproblems-problemphyxmini0601-deltawellscatteringsetup}
  Continuity of ψ and the derivative jump ψ'(0⁺)-ψ'(0⁻)=-(2mα/ℏ²)ψ(0) for every pair of incoming amplitudes. The outgoing amplitudes are obtained from the candidate physical S-matrix, so these laws determine that matrix without assuming its requested closed form
\end{definition}
\begin{proof}
  This is the declared model definition of Satisfies Delta Well Boundary Laws.
\end{proof}
\begin{definition}[Satisfies Scattering Matrix Role]
  \label{def:physics:phyx-mini-0601:phyxminiproblems-problemphyxmini0601-satisfiesscatteringmatrixrole}
  \lean{PhyXMiniProblems.ProblemPhyXMini0601.SatisfiesScatteringMatrixRole}
  \uses{def:physics:phyx-mini-0601:phyxminiproblems-problemphyxmini0601-deltawellscatteringsetup}
  The actual outgoing vector (B,F) equals S times the incoming (A,G)
\end{definition}
\begin{proof}
  This is the declared model definition of Satisfies Scattering Matrix Role.
\end{proof}
\begin{definition}[Satisfies Transfer Matrix Laws]
  \label{def:physics:phyx-mini-0601:phyxminiproblems-problemphyxmini0601-satisfiestransfermatrixlaws}
  \lean{PhyXMiniProblems.ProblemPhyXMini0601.SatisfiesTransferMatrixLaws}
  \uses{def:physics:phyx-mini-0601:phyxminiproblems-problemphyxmini0601-deltawellscatteringsetup}
  Transfer convention (F,G)ᵀ = M (A,B)ᵀ, together with the standard unit determinant law for a localized real one-dimensional potential
\end{definition}
\begin{proof}
  This is the declared model definition of Satisfies Transfer Matrix Laws.
\end{proof}
\begin{definition}[Is Left Incident Trial]
  \label{def:physics:phyx-mini-0601:phyxminiproblems-problemphyxmini0601-isleftincidenttrial}
  \lean{PhyXMiniProblems.ProblemPhyXMini0601.IsLeftIncidentTrial}
  \uses{def:physics:phyx-mini-0601:phyxminiproblems-problemphyxmini0601-deltawellscatteringsetup}
  The typical left-incident trial: G = 0 and the incident wave is nonzero
\end{definition}
\begin{proof}
  This is the declared model definition of Is Left Incident Trial.
\end{proof}
\begin{definition}[Satisfies Scattering Observables]
  \label{def:physics:phyx-mini-0601:phyxminiproblems-problemphyxmini0601-satisfiesscatteringobservables}
  \lean{PhyXMiniProblems.ProblemPhyXMini0601.SatisfiesScatteringObservables}
  \uses{def:physics:phyx-mini-0601:phyxminiproblems-problemphyxmini0601-deltawellscatteringsetup}
  Definitions of the four reflection/transmission coefficients from the source, plus the amplitude-ratio definition for the actual left-incident trial
\end{definition}
\begin{proof}
  This is the declared model definition of Satisfies Scattering Observables.
\end{proof}
\begin{definition}[Answer Choice]
  \label{def:physics:phyx-mini-0601:phyxminiproblems-problemphyxmini0601-answerchoice}
  \lean{PhyXMiniProblems.ProblemPhyXMini0601.AnswerChoice}
  Labels attached to the four expressions in the source
\end{definition}
\begin{proof}
  This is the declared model definition of Answer Choice.
\end{proof}
\begin{definition}[displayed Answer Value]
  \label{def:physics:phyx-mini-0601:phyxminiproblems-problemphyxmini0601-displayedanswervalue}
  \lean{PhyXMiniProblems.ProblemPhyXMini0601.displayedAnswerValue}
  \uses{def:physics:phyx-mini-0601:phyxminiproblems-problemphyxmini0601-deltawellscatteringsetup, def:physics:phyx-mini-0601:phyxminiproblems-problemphyxmini0601-answerchoice}
  The scalar expression printed beside each answer label
\end{definition}
\begin{proof}
  This is the declared model definition of displayed Answer Value.
\end{proof}
\begin{definition}[Matches Answer Choice]
  \label{def:physics:phyx-mini-0601:phyxminiproblems-problemphyxmini0601-matchesanswerchoice}
  \lean{PhyXMiniProblems.ProblemPhyXMini0601.MatchesAnswerChoice}
  \uses{def:physics:phyx-mini-0601:phyxminiproblems-problemphyxmini0601-deltawellscatteringsetup, def:physics:phyx-mini-0601:phyxminiproblems-problemphyxmini0601-answerchoice, def:physics:phyx-mini-0601:phyxminiproblems-problemphyxmini0601-displayedanswervalue}
  The modeled left-transmission coefficient agrees with a displayed choice
\end{definition}
\begin{proof}
  This is the declared model definition of Matches Answer Choice.
\end{proof}
\begin{definition}[recorded Dataset Answer]
  \label{def:physics:phyx-mini-0601:phyxminiproblems-problemphyxmini0601-recordeddatasetanswer}
  \lean{PhyXMiniProblems.ProblemPhyXMini0601.recordedDatasetAnswer}
  \uses{def:physics:phyx-mini-0601:phyxminiproblems-problemphyxmini0601-answerchoice}
  Recorded dataset metadata; it is not a physical-law premise
\end{definition}
\begin{proof}
  This is the declared model definition of recorded Dataset Answer.
\end{proof}
\begin{lemma}[attractive Delta Well scattering Matrix]
  \label{lem:physics:phyx-mini-0601:phyxminiproblems-problemphyxmini0601-attractivedeltawell-scatteringmatrix}
  \lean{PhyXMiniProblems.ProblemPhyXMini0601.attractiveDeltaWell_scatteringMatrix}
  \uses{def:physics:phyx-mini-0601:phyxminiproblems-problemphyxmini0601-deltawellscatteringsetup, def:physics:phyx-mini-0601:phyxminiproblems-problemphyxmini0601-hasphysicaldeltawellparameters, def:physics:phyx-mini-0601:phyxminiproblems-problemphyxmini0601-deltawellcoupling, def:physics:phyx-mini-0601:phyxminiproblems-problemphyxmini0601-satisfiesdeltawellboundarylaws}
  The delta boundary conditions determine the scattering matrix of the attractive potential -α δ(x). This is the construction requested by the written question, and its closed form occurs only in the conclusion
\end{lemma}
\begin{proof}
  The conclusion follows from the governing relations and figure readouts named in the statement of attractive Delta Well scattering Matrix.
\end{proof}
\begin{lemma}[left Transmission transfer Matrix Formula]
  \label{lem:physics:phyx-mini-0601:phyxminiproblems-problemphyxmini0601-lefttransmission-transfermatrixformula}
  \lean{PhyXMiniProblems.ProblemPhyXMini0601.leftTransmission_transferMatrixFormula}
  \uses{def:physics:phyx-mini-0601:phyxminiproblems-problemphyxmini0601-deltawellscatteringsetup, def:physics:phyx-mini-0601:phyxminiproblems-problemphyxmini0601-satisfiestransfermatrixlaws, def:physics:phyx-mini-0601:phyxminiproblems-problemphyxmini0601-isleftincidenttrial, def:physics:phyx-mini-0601:phyxminiproblems-problemphyxmini0601-satisfiesscatteringobservables}
  With the convention (F,G)ᵀ = M(A,B)ᵀ, left incidence and det M = 1 give the recorded transmission formula. Neither this formula nor answer C is a field of the transfer or observable laws
\end{lemma}
\begin{proof}
  The conclusion follows from the governing relations and figure readouts named in the statement of left Transmission transfer Matrix Formula.
\end{proof}
% --- Archon declaration topology end ---
% --- Archon physics formalization source end ---
